\section{Criterios de calidad del estudio}
\label{sec:calidad}

El rigor científico y metodológico del estudio se verifica mediante cuatro
criterios, cada uno con su mecanismo concreto de control y su evidencia de
verificación.

\subsection{Validez interna y trazabilidad}

\textbf{Definición.} La capacidad de atribuir las diferencias observadas en el
error exclusivamente al módulo evaluado, y de que un tercero obtenga exactamente
las mismas salidas.

\textbf{Mecanismo.} Aislamiento del entorno de ejecución, fijación determinista
de la semilla del generador pseudoaleatorio, derivación de flujos aleatorios
independientes por corrida a partir de una única semilla maestra, y
parametrización centralizada en una estructura inmutable.

\textbf{Verificación.} Se ejecutó el experimento completo dos veces de forma
independiente y se compararon los cuatro archivos CSV generados mediante suma de
verificación MD5, obteniendo identidad byte a byte. Este criterio, que en las
entregas previas se declaraba sin demostración, se encuentra ahora comprobado.

\subsection{Confiabilidad y convergencia asintótica}

\textbf{Definición.} La estabilidad de las conclusiones frente a la
variabilidad estocástica inherente al proceso.

\textbf{Mecanismo.} Diseño de Monte Carlo con $M = 50$ corridas independientes
de $N = 1000$ pasos, que asegura la convergencia de la ganancia
$\mathbf{K}_{k}$ y de la covarianza $\mathbf{P}_{k|k}$ hacia sus valores de
régimen permanente según la formulación de Riccati.

\textbf{Verificación.} La figura~\ref{fig:convergencia} documenta la
convergencia efectiva de ambas magnitudes. La dispersión de las métricas entre
corridas es reducida —desviación estándar del RMSE de \SI{0,059}{\metre\cubed}
para el filtro— y la comparación pareada de la figura~\ref{fig:montecarlo}
muestra que la ventaja del módulo se sostiene en las cincuenta corridas sin
excepción.

\subsection{Validez de constructo}

\textbf{Definición.} La correspondencia entre el modelo matemático implementado
y el fenómeno físico que pretende representar.

\textbf{Mecanismo.} Las matrices del modelo se derivan estrictamente de la ley
de conservación de masa de la tolva, y no de un ajuste empírico a los datos. Los
parámetros dimensionales se calibran contra la escala real del activo: la tasa
de extracción corresponde a la capacidad documentada de un chancador giratorio
de gran minería, y el volumen por evento de descarga al tonelaje de un camión de
extracción, según se justifica en la tabla~\ref{tab:parametros}.

\textbf{Verificación.} Además de la trazabilidad física de cada parámetro, se
verifica computacionalmente la condición de observabilidad del par
$(\mathbf{A}, \mathbf{H})$ antes de cada corrida
(ecuación~\eqref{eq:observabilidad}), y se comprobó que el MSE de la telemetría
cruda coincide con la varianza de ruido inyectada, lo que confirma la
implementación correcta del generador de degradación.

Corresponde señalar aquí, como control de calidad efectivamente aplicado, que
este criterio condujo a la detección y corrección de la inconsistencia física de
la covarianza de proceso documentada en la sección~\ref{sec:recalibracion}.

\subsection{Sensibilidad y robustez}

\textbf{Definición.} Los márgenes dentro de los cuales el módulo conserva su
utilidad frente a variaciones del entorno y errores de calibración.

\textbf{Mecanismo.} Dos experimentos independientes: barrido de la varianza real
del sensor sobre dos órdenes de magnitud, y malla de desajuste intencional de
las razones de covarianza $\mathbf{Q}/\mathbf{R}$ sobre un rango de cuatro
órdenes de magnitud.

\textbf{Verificación.} Los resultados de las tablas~\ref{tab:degradacion} y de
la figura~\ref{fig:sensibilidad} acreditan que el aporte del módulo crece con la
degradación del entorno y que veintidós de veinticinco combinaciones de
sintonización superan a la línea base. Este criterio, declarado en el diseño
metodológico de la Fase 3, se implementa efectivamente en esta entrega.

\subsection{Limitaciones de los criterios de calidad}

La honestidad metodológica exige declarar el alcance de estos criterios. La
validez interna del estudio es alta, pero su \textbf{validez externa} está
acotada por construcción: los resultados acreditan el desempeño del estimador
sobre un proceso que satisface exactamente los supuestos del modelo —linealidad,
ruido gaussiano blanco, covarianzas conocidas—, condición que una planta real
solo cumple de manera aproximada. El análisis de robustez de la
sección~\ref{sec:resultados} mitiga parcialmente esta limitación al evaluar el
desempeño bajo desajuste de covarianzas, pero no sustituye la validación con
telemetría real, que queda declarada como trabajo futuro en la
sección~\ref{sec:conclusiones}.
